Assume all displayed hypotheses and let
\[
 \mathcal B=\{A\in\pi_R:A\subseteq B\}.
\]
Because \(\pi_R\preceq\pi_{R'}\), the blocks in \(\mathcal B\) partition
\(B\).  By hypothesis \(|\mathcal B|\geq2\).  The subgraph of the finite
acyclic graph \(G_R\) induced by \(\mathcal B\) is a finite DAG, hence has a
sink; choose one and denote it by \(D\).  Put \(C=B\setminus D\).  Then
\(B=C\mathbin{\dot\cup}D\) is a nontrivial coordinate split because \(D\) is
a nonempty partition block and at least one other block lies in \(C\).

Every parent of a fine block in \(B\) lies in \(B\) by fine-parent-closure.
Moreover, \(D\) is not a parent of any block in \(\mathcal B\setminus\{D\}\),
because \(D\) is a sink of the induced graph.  Consequently every parent of
every fine block contained in \(C\) is itself contained in \(C\).  Recursing
through a topological ordering of this induced parent-closed subgraph and using
the structural equations shows that
\[
 X_C\quad\hbox{is measurable with respect to}\quad
 \sigma\!\left(\varepsilon_{R,A}:A\in\mathcal B, A\subseteq C\right).
\]
The same parent-closure hypothesis gives
\(\bigcup\operatorname{pa}_{G_R}(D)\subseteq C\).  By mutual independence of
the \(R\)-block innovations,
\(\varepsilon_{R,D}\) is independent of the sigma-field in the last display,
and therefore
\[
 \varepsilon_{R,D}\perp\!\!\!\perp X_C. \tag{1}
\]

Since \(B\) is a root of \(G_{R'}\), its mechanism has a zero-dimensional
domain and equals a constant vector \(c_B\in\mathbb R^{|B|}\).  Write
\(c_C\) and \(c_D\) for its coordinate subvectors.  The coarse root equation
is
\[
 X_B=c_B+\varepsilon_{R',B},
\]
so
\[
 \varepsilon_{R',C}=X_C-c_C,
 \qquad
 \varepsilon_{R',D}=X_D-c_D. \tag{2}
\]
The fine equation for \(D\) and the parent containment just proved give
\[
 X_D=f_{R,D}\!\left(X_{\cup\operatorname{pa}_{G_R}(D)}\right)
       +\varepsilon_{R,D}. \tag{3}
\]
For \(u_C\in\mathbb R^{|C|}\), define
\[
 h(u_C)=
 f_{R,D}\!\left((u_C+c_C)_{\cup\operatorname{pa}_{G_R}(D)}\right)-c_D.
\]
The coordinate projection and translation are smooth, and
\(f_{R,D}\in C^3\) by the admissible-class mechanism assumption; hence
\(h\in C^3(\mathbb R^{|C|},\mathbb R^{|D|})\).  Substituting (2) into (3)
yields the asserted identity
\[
 \varepsilon_{R',D}
 =h(\varepsilon_{R',C})+\varepsilon_{R,D}. \tag{4}
\]
Equation (1) and the first equality in (2) imply
\[
 \varepsilon_{R,D}\perp\!\!\!\perp\varepsilon_{R',C}. \tag{5}
\]

It remains to check the residual-law quantifier.  Since \(R\in\mathfrak A(P)\),
the law of \(\varepsilon_{R,D}\) has a unit-integral Lebesgue density that is
strictly positive on \(\mathbb R^{|D|}\), and its log-density belongs to
\(C^3(\mathbb R^{|D|})\).  By
Definition~\ref{def:smooth-positive-laws},
\[
 \mathcal L(\varepsilon_{R,D})\in\operatorname{SmoothPos}(|D|). \tag{6}
\]
Equations (4)--(6) form a directional additive-noise representation for the
nontrivial split \(B=C\mathbin{\dot\cup}D\) of the coarse innovation
\(\varepsilon_{R',B}\).  It is forbidden by innovation atomicity of \(R'\).
The resulting contradiction proves the theorem.  The root equation was used
in (2), and fine-parent-closure was used both before (1) and before (3); the
argument therefore makes no assertion when either condition is absent.